\documentclass[10pt,a4paper]{amsart}
\usepackage[margin=19mm]{geometry}
\usepackage{amsmath,amssymb,mathtools}
\usepackage[scr=rsfs,cal=euler]{mathalfa}
\usepackage{xfrac}
\usepackage[unicode,hidelinks,bookmarks=false]{hyperref}
\newcommand{\ie}{i.e.\ }
\newcommand{\cf}{cf.\ }
\newcommand{\resp}{respectively\ }
\newcommand{\Subsection}{\subsection}
\providecommand{\nobreakdash}{\nobreak\hbox{-}}
\setlength{\emergencystretch}{2em}
\multlinegap0pt
\usepackage{amssymb}
\usepackage{siunitx}
\newtheorem{theorem}{Theorem}
\title{Some general results on risk budgeting portfolios}
\author{Claudia Fassino, Pierpaolo Uberti}
\date{Source: arXiv:2603.15511v1 (CC BY 4.0)}
\begin{document}
\maketitle

\noindent\emph{Source and licence.} Some general results on risk budgeting portfolios. Version arXiv:2603.15511v1 (CC BY 4.0), distributed by arXiv under the Creative Commons Attribution 4.0 International licence, http://creativecommons.org/licenses/by/4.0/, as recorded in the arXiv licence metadata for this exact version and shown on the abstract page https://arxiv.org/abs/2603.15511v1. Full licence notice: \texttt{licenses/arxiv-2603.15511-CC-BY-4.0.txt}.\par\smallskip

\noindent\textit{Editorial scope.} Counted unit: the article's theorem normrest (source0.tex lines 571--581). The article's complete original proof of it is delivered with it (source0.tex lines 582--621, one \texttt{proof} environment of 40 lines). No mathematics is written, repaired or re-derived: the statement and the proof are the article's own lines, in the article's own notation, and no retained line is substituted or re-flowed.  No further source-local context is needed: every cross-reference of every retained line resolves inside the two delivered spans.  Nothing was omitted from any retained span, so the package carries no omission record.  No retained line contains a citation, so the wrapper carries no bibliography and no bibliography entry.  No retained line contains a figure environment, an image-inclusion command or drawing code, so no image is delivered and the package declares no asset dependency.  Editorial formatting only, in addition: an AMS \texttt{amsart} preamble that restates the article's own declarations and macros used by the retained lines, namely the environments \texttt{theorem} on separate counters, together with the standard packages those lines need.  The retained environments print in this document as Theorem 1 in order of appearance, whereas the article prints the same items with the numbering of its own full document.  The complete native source of the article as distributed in the arXiv e-print for this exact version is bundled unchanged as \texttt{sources/196484/source0.tex}.
\begin{theorem}\label{normrest}
If $\mathbf{x}_n \in S$, then $\mathbf{x}_{n+1}=\mathbf{x}_n+k_n\Delta(\mathbf{x}_n) \in S$ if 
%Sia $x_n$ appartenente al simplesso. Allora $x_{n+1}=x_n+k_n\Delta(x_n)$ appartiene al simplesso se
\begin{eqnarray*}
|k_n| \le \left\{\begin{array}{l} \min \left \{ \frac{1-(\mathbf{x}_n)_i}{ |\Delta(\mathbf{x}_n)_i | }  \ | \ k_n\Delta(\mathbf{x}_n)_i >0  \right  \} \\ \\
 \min \left \{ \frac{(\mathbf{x}_n)_i}{ |\Delta(\mathbf{x}_n)_i| }  \ | \ k_n\Delta(\mathbf{x}_n)_i <0  \right  \} \ ,
\end{array}
\right.
\end{eqnarray*}
where $(\mathbf{x}_n)_i$ and $\Delta(\mathbf{x}_n)_i$ are the $i^{th}$-entries of $\mathbf{x}_n$ and $\Delta(\mathbf{x}_n)$ respectively.
\end{theorem}

\begin{proof}
First,
$ \mathbf 1^t \Delta(\mathbf{x}_n)=\mathbf{x}_n^tV\mathbf{x}_n-(\mathbf 1^t \mathbf{x}_b)(\mathbf{x}_n^tV\mathbf{x}_n)=\mathbf{x}_n^tV\mathbf{x}_n-\mathbf{x}_n^tV\mathbf{x}_n=0$; this means that if the entries of $\mathbf{x}_n$ have unitary sum, $\mathbf{x}_{n+1}$ has the same property if $\mathbf 1^t \mathbf{x}_{n+1}=\mathbf 1^t(\mathbf{x}_n+k_n \Delta(\mathbf{x}_n))=1$, 
independently from the value of $k_n$. Therefore, it is necessary to prove that each component of $\mathbf{x}_{n+1}$ is non-negative. For simplicity, we omit the index $n$ and refer to $\Delta$ for $\Delta(\mathbf{x})$. We require 
$ 0 \le \mathbf{x}_i +k \Delta_i \le 1$, i.e. $ -\mathbf{x}_1 \le k \Delta_i \le 1-\mathbf{x}_i$, 
where $\mathbf{x}_i$ and $\Delta_i$ are the $i^{th}$ entries of $\mathbf{x}$ and $\Delta$ respectively. By assumption, $\mathbf{0} < \mathbf{x}_1 < \mathbf{1}$.
Two cases are possible.
\smallskip\noindent
\textbf{case $1$}: $k \Delta_i >0$, then $k \Delta_i =|k \Delta_i |= |k|| \Delta_i |$. The value of $|k|$ needs to satisfy the conditions 
\begin{eqnarray*}
\left\{\begin{array}{rcl} |k|&\ge& -\frac{\mathbf{x}_i}{ |\Delta_i| } \\ \\
|k| & \le & \frac{1-\mathbf{x}_i}{ |\Delta_i| }
\end{array}
\right. 
\end{eqnarray*}
Note that the first inequality is always verified.

\smallskip\noindent
\textbf{case $2$}: $k \Delta_i <0$, then $k \Delta_i =-|k \Delta_i |= -|k|| \Delta_i |$. The value of $|k|$ needs to satisfy the conditions 
\begin{eqnarray*}
\left\{\begin{array}{rcl} -|k|&\ge& -\frac{\mathbf{x}_i}{ |\Delta_i| } \\ \\
-|k| & \le & \frac{1-\mathbf{x}_i}{ |\Delta_i| }
\end{array}
\right. \qquad \Leftrightarrow \qquad  \left\{\begin{array}{rcl} |k|&\le& \frac{\mathbf{x}_i}{ |\Delta_i| } \\ \\
|k| & \ge & \frac{\mathbf{x}_i-1 }{ |\Delta_i| }
\end{array}
\right. 
\end{eqnarray*}
In this case, the second inequality is always verified.


\smallskip\noindent
Concluding, $|k|$ verifies 
\begin{eqnarray*}
|k| \le \left\{\begin{array}{l} \min \left \{ \frac{1-\mathbf{x}_i}{ |\Delta_i| }  \ | \ k\Delta_i >0  \right  \} \\ \\
 \min \left \{ \frac{\mathbf{x}_i}{ |\Delta_i| }  \ | \ k\Delta_i <0  \right  \}
\end{array}
\right.
\end{eqnarray*}
\end{proof}

\end{document}
